\documentclass{article}
\usepackage{geometry}
\geometry{a4paper, margin=1in}
\usepackage{hyperref}
\hypersetup{colorlinks=true, urlcolor=blue, citecolor=blue}

\title{Proposal B: Spatiotemporal Masked Autoencoding for\\Embryo Time-Lapse Microscopy}
\author{Candidate Name}
\date{Suggested Supervisor: Prof.\ Dr.\ Martin Weigert}

\begin{document}
\maketitle

\section*{Motivation}

Time-lapse microscopy of embryonic development produces massive video datasets where the temporal dimension carries as much biological information as the spatial dimensions. Yet self-supervised pretraining for microscopy has focused almost exclusively on \textbf{spatial} objectives: masked autoencoders (MAE) reconstruct masked \emph{image patches}, contrastive methods compare \emph{static augmentations}, and even TAP~\cite{tap} predicts temporal order from \emph{individual frames}. None of these methods learn to predict \textbf{across time}---reconstructing future or masked frames from temporal context.

In natural video understanding, \textbf{VideoMAE}~\cite{videomae} demonstrated that masking and reconstructing spatiotemporal tubes in video produces powerful representations for action recognition and temporal reasoning. No one has applied this idea to microscopy time-lapse data. The opportunity is clear: embryonic development has strong temporal structure (cell division, migration, tissue deformation) that a spatiotemporal pretext task should capture.

The Weigert group recently introduced \textbf{MuViT}~\cite{muvit}, a multi-resolution Vision Transformer with MAE pretraining for gigapixel microscopy images (CVPR 2026). MuViT uses RoPE positional embeddings and a modular ViT encoder---both of which naturally extend to a temporal dimension. This project builds directly on MuViT's codebase to create a \textbf{spatiotemporal MAE} for microscopy video.

\section*{Approach}

\begin{enumerate}
	\item \textbf{Spatiotemporal patch embedding.} Extend MuViT's patch embedding to include a temporal dimension. A video clip of $T$ frames is divided into spatiotemporal tubes (e.g., $16 \times 16 \times 4$ patches covering space and time). Each tube is projected to a token.
	\item \textbf{Temporal RoPE extension.} MuViT uses Rotary Position Embeddings (RoPE) with world coordinates for multi-resolution spatial alignment. Extend these to encode temporal position, enabling the transformer to reason about both spatial location and time.
	\item \textbf{Spatiotemporal masking and reconstruction.} Mask a high proportion (e.g., 90\%) of spatiotemporal tubes---following VideoMAE's insight that high masking ratios force temporal reasoning rather than spatial interpolation. The encoder processes only visible tubes; a lightweight decoder reconstructs masked tubes.
	\item \textbf{Pretraining on SLICE-1 embryo data.} Train on unlabeled time-lapse clips from the SLICE-1 dataset~\cite{slice1} (50 \textit{Tribolium} embryos). No annotations needed.
	\item \textbf{Downstream evaluation.} Evaluate learned representations on:
	      \begin{itemize}
		      \item \textbf{Cell tracking:} Extract per-cell features and feed into Trackastra~\cite{trackastra}; compare against TAP and SAM2 baselines.
		      \item \textbf{Developmental stage classification:} Linear probe on the CLS token to classify known \textit{Tribolium} stages (blastoderm, gastrulation, germband extension, etc.).
		      \item \textbf{Temporal prediction quality:} Qualitatively inspect reconstructed masked future frames---does the model learn biologically plausible dynamics?
	      \end{itemize}
\end{enumerate}

\section*{Why This Works as a One-Semester Project}

\begin{itemize}
	\item \textbf{MuViT provides the starting codebase}: the ViT encoder, MAE training loop, RoPE implementation, and microscopy data loading are already implemented in PyTorch. The core engineering is extending spatial MAE to spatiotemporal MAE---a well-understood architectural modification.
	\item \textbf{VideoMAE is a mature reference}: the spatiotemporal masking strategy, optimal masking ratios, and tube design are documented in the natural-video literature. The novelty is the application domain, not the method itself.
	\item \textbf{SLICE-1 provides abundant data}: 5.7\,M images across 50 embryos. Pretraining data is not a bottleneck.
	\item The project is \textbf{incrementally buildable}: (1)~reproduce MuViT's spatial MAE on microscopy frames, (2)~extend to spatiotemporal, (3)~pretrain, (4)~evaluate downstream.
\end{itemize}

\section*{Expected Contribution}

The first application of masked spatiotemporal autoencoding to microscopy video. This directly tests whether \textbf{predicting across time} produces more biologically meaningful representations than spatial-only pretraining (MAE) or temporal-order prediction (TAP). If the spatiotemporal MAE learns to reconstruct cell divisions and tissue deformations, it demonstrates that embryo dynamics can be captured by self-supervision alone---opening the door to foundation models for developmental biology.

\begin{thebibliography}{9}
	\bibitem{tap} B.~Gallusser, M.~Stieber, and M.~Weigert. ``Self-supervised dense representation learning for live-cell microscopy with time arrow prediction.'' \textit{MICCAI}, 2023.
	\bibitem{videomae} Z.~Tong et al. ``VideoMAE: Masked Autoencoders are Data-Efficient Learners for Self-Supervised Video Pre-Training.'' \textit{NeurIPS}, 2022.
	\bibitem{muvit} A.~Dominguez Mantes, G.~La Manno, and M.~Weigert. ``MuViT: Multi-Resolution Vision Transformer for Gigapixel Microscopy.'' \textit{CVPR}, 2026.
	\bibitem{trackastra} B.~Gallusser and M.~Weigert. ``Trackastra: Transformer-based cell tracking for live-cell microscopy.'' \textit{ECCV}, 2024.
	\bibitem{slice1} F.~Strobl et al. ``SLICE-1: A comprehensive light-sheet imaging dataset of \textit{Tribolium castaneum} embryogenesis.'' \textit{Scientific Data}, Dec 2025.
\end{thebibliography}

\end{document}
